\chapter{Architectural Lineage of CSDI}
\label{app:arch_lineage}

Section~\ref{sec:csdi_literature} described the four-step lineage PixelCNN, WaveNet, DiffWave, CSDI. Here the explanation is expanded, focusing on which ideas each introduced and which it
passed forward.

\subsection{PixelCNN}

PixelCNN \cite{oord2016pixel} generates images autoregressively by factorising
the joint over pixels in Raster order (left to right, row by row, top to bottom). To implement this with a convolutional
network, it uses \textit{masked convolutions}: the convolutional kernel is
multiplied by a binary mask that zeroes out weights corresponding to future
pixels, ensuring the receptive field never sees what it should be predicting.\\

Two components introduced by PixelCNN are inherited by every subsequent model:
\begin{itemize}
    \item \textit{Gated activations}: $\tanh(W_f * h) \odot \sigma(W_g * h)$,
    where $\sigma(\cdot)$ is a selection gate $\in [0,1]$ and $\tanh(\cdot)$
    is a content filter $\in [-1,1]$. Their product suppresses uninformative
    activations, which is especially useful at low signal-to-noise ratio.
    \item \textit{Residual connections}: each block learns a residual correction
    on top of its input, enabling deep stacks without vanishing gradients.
\end{itemize}
PixelCNN's limitation for long sequences is its sequential inference: each new
token requires conditioning on all previous ones, giving $O(L)$ passes through
the network for a sequence of length $L$.

\subsection{WaveNet}

WaveNet \cite{oord2016wavenet} transfers PixelCNN's gated residual structure
to raw audio by replacing 2D masked image convolutions with 1D
\textit{causal dilated convolutions}. A dilated convolution with dilation $d$
skips $d-1$ samples between each kernel element. Stacking dilations
$d = 1, 2, 4, \ldots, 2^{L-1}$ within a residual block gives an exponentially
growing receptive field of $2^L$ samples using only $O(L)$ parameters per
block, while causality is preserved by ensuring the kernel never reaches
forward in time.

\subsection{DiffWave}

DiffWave \cite{kong2021diffwave} adopts the WaveNet-style residual blocks
(gated activations, dilated convolutions, skip aggregation, conditioning
injection) but removes the causality constraint. Since a diffusion model
denoises the entire signal at once, not sequentially, future context is
available and causal convolutions become unnecessary. Replacing them with
\textit{bidirectional} dilated convolutions allows the receptive field to
look both forward and backward, capturing symmetric temporal context.\\

The diffusion step index $t$ replaces the speaker-identity conditioning: it
is sinusoidally embedded and injected into each residual block, telling the
network its current noise level (identical in function to the diffusion-step
embedding in CSDI, Section~\ref{sec:csdi_literature}). DiffWave is the direct
predecessor of CSDI's backbone: CSDI replaces DiffWave's 1D bidirectional
dilated convolutions with the 2D factored attention (temporal $\times$
feature axes), and adds the conditional masking mechanism, as described in
Section~\ref{sec:csdi_literature}.